\section{Main Theorem}\label{sec:main-theorem}

\begin{theorem}[Anchored derivative-closure coefficient sweep]\label{thm:main}
The static and homogeneous clause below holds under
Assumptions~\ref{assump:parameter-regime}--\ref{assump:anchored-derivative-closure}.
All law-bearing clauses also use Assumption~\ref{assump:cube-density-laws}.
In every case the deterministic presentation is fixed before the law and interval are selected.

\emph{Static and homogeneous certificate.} For every $\theta\in\Theta$,
\begin{equation}
\label{eq:main-display-001}
\|B(\theta)\|_{\mathrm{op}}\leq\|B(\theta)\|_{\mathrm F}\leq\widehat\Lambda_{B,T},
\qquad
\sup_{\theta\in\Theta}\|B(\theta)\|_{\mathrm{op}}\leq\widehat\Lambda_{B,T}.
\end{equation}
If $F_0\equiv0$ and $B_F$ is the block indexed by $1,\ldots,N$, then
\begin{equation}
\label{eq:main-display-002}
F'=B_FF,\qquad
\gamma_F'=(I_N-\gamma_F\gamma_F^{\mathsf T})B_F\gamma_F,\qquad
\Gamma_{\mathrm{proj}}(F)\leq\widehat\Lambda_{B,T}.
\end{equation}
The anchor gives $F_{j_*}=1$, $\|F\|_2\geq1$, and $F\neq0$.

\emph{Equivalent affine chart form.} For every $\mu\in\mathcal D$ and every interval
$I\subseteq\Theta$ with $|I|>0$, and every Lebesgue-measurable partition
$I=\bigsqcup_{j=1}^NE_j$ with $F_j\neq0$ on $E_j$,
\begin{equation}
\label{eq:main-display-003}
\begin{aligned}
\Pr_{\alpha\sim\mu}[\exists\theta\in I:\phi_\alpha(\theta)=0]
&\leq\kappa\sum_{j=1}^N\int_{E_j}\int_{[-R,R]^{N-1}}
\mathbf 1\{|T_j(\theta,\beta)|\leq R\}|\partial_\theta T_j(\theta,\beta)|\,d\beta\,d\theta\\
&\leq\kappa\sum_{j=1}^N\int_{E_j}\int_{[-R,R]^{N-1}}
|\partial_\theta T_j(\theta,\beta)|\,d\beta\,d\theta.
\end{aligned}
\end{equation}
These are ordinary-probability inequalities in the original $N$ coefficient coordinates; they remain valid
with no pivot margin, including tangent, multiple, endpoint, empty-cell, zero-Jacobian, persistent-root, and
$N=1$ cases.

\emph{Coordinate-free affine rate.} For the same $\mu$ and $I$,
\begin{equation}
\label{eq:main-display-004}
\begin{aligned}
\Pr_{\alpha\sim\mu}[\exists\theta\in I:\phi_\alpha(\theta)=0]
&\leq\kappa\int_I\int_{H_\theta\cap[-R,R]^N}
\frac{|F_0'(\theta)+\langle a,F'(\theta)\rangle|}{\|F(\theta)\|_2}
\,d\mathcal H^{N-1}(a)\,d\theta\\
&\leq\kappa\sqrt2(2R)^{N-1}(1+NR^2)\widehat\Lambda_{B,T}|I|\\
&=\frac{A(1+NR^2)\widehat\Lambda_{B,T}}{\sqrt2R}|I|.
\end{aligned}
\end{equation}
Consequently,
\begin{equation}
\label{eq:main-display-005}
C^{\mathrm{aff}}_{\mathcal D}(F_0,F;\Theta)
\leq\frac{A(1+NR^2)\widehat\Lambda_{B,T}}{\sqrt2R}.
\end{equation}
If $\widehat\Lambda_{B,T}=0$, the event probability is exactly zero. The section convention is Euclidean
$\mathcal H^{N-1}$, with $\mathcal H^0$ counting measure when $N=1$.

\emph{Sharper homogeneous rate.} If $F_0\equiv0$, then for $a\in H_\theta$,
\begin{equation}
\label{eq:main-display-006}
H_\theta=F(\theta)^\perp=\gamma_F(\theta)^\perp,\qquad
\frac{|\langle a,F'(\theta)\rangle|}{\|F(\theta)\|_2}=|\langle a,\gamma_F'(\theta)\rangle|,
\end{equation}
and
\begin{equation}
\label{eq:main-display-007}
\Pr_{\alpha\sim\mu}[\exists\theta\in I:\langle\alpha,F(\theta)\rangle=0]
\leq A\sqrt{\frac N2}\Gamma_{\mathrm{proj}}(F)|I|
\leq A\sqrt{\frac N2}\widehat\Lambda_{B,T}|I|,
\end{equation}
\begin{equation}
\label{eq:main-display-008}
C^{\mathrm{Pf}}_{\mathcal D}(F;\Theta)
\leq A\sqrt{\frac N2}\widehat\Lambda_{B,T}.
\end{equation}
The interval supremum is taken for each fixed law before the law supremum. If
$\Gamma_{\mathrm{proj}}(F)=0$, every nonempty-interval event is one fixed proper central hyperplane and has
probability zero under every admissible full density.

\emph{Exact affine-monic specialization.} For every integer $d\geq1$, every bounded interval
$J\subset\mathbb R$, and every possibly correlated law on $[-R,R]^d$ with full density bounded by $\kappa$,
set
\begin{equation}
\label{eq:main-display-009}
F_0(\theta)=\theta^d,\qquad F_{k+1}(\theta)=\theta^k\ (0\leq k\leq d-1),\qquad
p_\alpha(\theta)=\theta^d+\sum_{k=0}^{d-1}\alpha_k\theta^k.
\end{equation}
Then $q=M=m=0$, $\Delta=N=d$, $A=(2R)^d\kappa$, the random vector is exactly
$\alpha=(\alpha_0,\ldots,\alpha_{d-1})$, and the leading coefficient one is deterministic and outside that
vector. The constant $(d+1)$-square shift has only
\begin{equation}
\label{eq:main-display-010}
B_{0,d}=d,\qquad B_{k+1,k}=k\quad(1\leq k\leq d-1),\qquad
\widehat\Lambda_{B,T}=\left(\sum_{k=1}^d k^2\right)^{1/2}.
\end{equation}
For $d\geq2$, use $E_1=J\cap\{|\theta|\leq1\}$ and
$E_d=J\cap\{|\theta|>1\}$ (all intermediate cells are empty), with
\begin{equation}
\label{eq:main-display-011}
T_1(\theta,\beta)=-\theta^d-\sum_{k=1}^{d-1}\beta_k\theta^k,
\qquad
T_d(\theta,\beta)=-\theta-\sum_{k=0}^{d-2}\beta_k\theta^{k-d+1}.
\end{equation}
Their velocities obey
\begin{equation}
\label{eq:main-display-012}
|\partial_\theta T_1|\leq d+\frac{Rd(d-1)}2\quad(|\theta|\leq1),
\end{equation}
\begin{equation}
\label{eq:main-display-013}
|\partial_\theta T_d|\leq1+R\sum_{m=1}^{d-1}\frac{m}{|\theta|^{m+1}}
\leq1+\frac{Rd(d-1)}2\leq d+\frac{Rd(d-1)}2\quad(|\theta|>1).
\end{equation}
At $d=1$, $[-R,R]^0$ has mass one, $T_1(\theta)=-\theta$, and
$|\partial_\theta T_1|=1$. In every case, including empty and singleton $J$,
\begin{equation}
\label{eq:main-display-014}
\Pr_{\alpha\sim\mu}[\exists\theta\in J:p_\alpha(\theta)=0]
\leq\kappa(2R)^{d-1}\left(d+\frac{Rd(d-1)}2\right)|J|.
\end{equation}

\emph{Counter-example 1.} For $0<\delta\leq1$, take $\Theta=[-1,1]$, $F_0=0$,
$F=(1,\theta/\delta)$, $R=1$, $\kappa=1/4$, and $A=1$. Then
\begin{equation}
\label{eq:main-display-015}
q=M=0,\quad \Delta=1,\quad N=2,\quad m=0,
\qquad \widehat\Lambda_{B,T}=\Gamma_{\mathrm{proj}}(F)=\frac1\delta,
\end{equation}
with the sole nonzero shift entry $B_{2,1}=1/\delta$. Every possibly correlated admissible law and every
positive-length interval $I\subseteq[-1,1]$ obey
\begin{equation}
\label{eq:main-display-016}
\Pr_{\alpha\sim\mu}[\exists\theta\in I:\alpha_1+\alpha_2\theta/\delta=0]\leq\frac{|I|}{\delta},
\qquad C^{\mathrm{Pf}}_{\mathcal D}(F;[-1,1])\leq\frac1\delta.
\end{equation}
For the selected independent uniform law on $[-1,1]^2$ and $0<\epsilon\leq\delta$,
\begin{equation}
\label{eq:main-display-017}
\Pr[\exists\theta\in[0,\epsilon]:\alpha_1+\alpha_2\theta/\delta=0]
=\frac{\epsilon}{4\delta},\qquad
\frac{\Pr[\exists\theta\in[0,\epsilon]:\alpha_1+\alpha_2\theta/\delta=0]}{\epsilon}
=\frac1{4\delta}.
\end{equation}
The geometric Ball-section normalization combined with the same certificate is the separate scale
$\sqrt2/\delta$; it is a section-geometry scale, not the all-law probability coefficient or the selected-law
ratio. Thus $1/(4\delta)$, $1/\delta$, and $\sqrt2/\delta$ are not identified with one another.

All assertions use ordinary probability for each fixed law, arbitrary full-joint coefficient correlation,
and literal endpoint conventions. Vector, operator, Frobenius, and scalar norms are Euclidean; the measures
are Lebesgue and Euclidean Hausdorff measures. No confidence parameter, hidden constant, chart-count factor,
interval-location factor, auxiliary tolerance, random leading coordinate, or second probability conversion is
present. Once the supplied certificate is fixed, additional dependence on $q,M,\Delta$ is degree zero. The result preserves both
baseline-invariance obligations: the exact deterministic-leading-coefficient monic theorem and the distinct
Counter-example scales. It makes no claim that an unrestricted raw Pfaffian presentation admits the supplied
derivative-closure certificate or a polynomial presentation-format bound.
\end{theorem}
